\documentclass[sigconf]{acmart}

%% DAC 2026 Conference Metadata
\setcopyright{acmcopyright}
\acmConference[DAC '26]{Proceedings of the 63rd Annual ACM/IEEE Design Automation Conference}{July 26--29, 2026}{Long Beach, CA, USA}
\acmBooktitle{Proceedings of the 63rd Annual ACM/IEEE Design Automation Conference (DAC '26)}
\acmYear{2026}
\acmDOI{}
\acmISBN{}

%%
%% \BibTeX command to typeset BibTeX logo in the docs
\AtBeginDocument{%
  \providecommand\BibTeX{{%
    Bib\TeX}}}

\input{./files/util.tex}
\renewcommand{\figurename}{Fig.}
\renewcommand{\tablename}{Tab.}
\usepackage{subcaption}
\setlength{\textfloatsep}{8pt plus 1pt minus 2pt}
\setlength{\floatsep}{6pt plus 1pt minus 2pt}
\setlength{\intextsep}{8pt plus 1pt minus 2pt}

\begin{document}

\title{GPU-Fuzz: Finding Memory Errors in Deep Learning Frameworks}

\author{Zihao Li$^{1,2,*}$, Hongyi Lu$^{1,2,3,*}$, Yanan Guo$^{4}$, Zhenkai Zhang$^{5}$, Shuai Wang$^{3}$, Fengwei Zhang$^{2,1,\dagger}$}
\email{{lizh2025, luhy2017}@mail.sustech.edu.cn}
\email{{yanan.guo}@rochester.edu, {zhenkai}@clemson.edu}
\email{{shuaiw}@cse.ust.hk, {zhangfw}@sustech.edu.cn}
\affiliation{%
  $^1$\textit{Research Institute of Trustworthy Autonomous Systems, Southern University of Science and Technology} \\
  $^2$\textit{Department of Computer Science and Engineering, Southern University of Science and Technology} \\
  $^3$\textit{Department of Computer Science and Engineering, Hong Kong University of Science and Technology} \\
  $^4$\textit{University of Rochester} \\
  $^5$\textit{Clemson University} \\
  \country{}
}

\begin{abstract}
  GPU memory errors are a critical threat to deep learning (DL) frameworks, leading to crashes or even security issues. We introduce \textsc{GPU-Fuzz}, a fuzzer locating these issues efficiently by modeling operator parameters as formal constraints. \textsc{GPU-Fuzz} utilizes a constraint solver to generate test cases that systematically probe error-prone boundary conditions in GPU kernels. Applied to PyTorch, TensorFlow, and PaddlePaddle, we uncovered 13 unknown bugs, demonstrating the effectiveness of \textsc{GPU-Fuzz} in finding memory errors.
  \end{abstract}

\keywords{GPU, Fuzzing, Deep Learning}

\begin{CCSXML}
<ccs2012>
 <concept>
  <concept_id>10010147.10010257.10010293.10010294</concept_id>
  <concept_desc>Computing methodologies~Machine learning</concept_desc>
  <concept_significance>500</concept_significance>
 </concept>
 <concept>
  <concept_id>10010147.10010257</concept_id>
  <concept_desc>Computing methodologies~Artificial intelligence</concept_desc>
  <concept_significance>500</concept_significance>
 </concept>
 <concept>
  <concept_id>10011007.10011006.10011008.10011009</concept_id>
  <concept_desc>Software and its engineering~Software safety</concept_desc>
  <concept_significance>500</concept_significance>
 </concept>
</ccs2012>
\end{CCSXML}

\ccsdesc[500]{Computing methodologies~Machine learning}
\ccsdesc[500]{Computing methodologies~Artificial intelligence}
\ccsdesc[500]{Software and its engineering~Software safety}

\maketitle

\begingroup
\renewcommand\thefootnote{}
\footnotetext{$^*$Zihao Li and Hongyi Lu contributed equally to this work.\\
$^\dagger$Fengwei Zhang is the corresponding author.}
\endgroup


\input{./files/1_intro.tex}
\input{./files/2_bg.tex}
\input{./files/3_design.tex}
\input{./files/4_impl.tex}
\input{./files/5_eval.tex}
\input{./files/6_discussion.tex}
\input{./files/7_related_work.tex}
\input{./files/8_disclosure.tex}
\input{./files/9_conclusion.tex}

\begin{acks}
We would like to thank the COMPASS members for their insightful comments. This work is partly supported by the National Natural Science Foundation of China under Grant No. U2541211, No. 62372218, and No. U24A6009. This work was also supported in part by a grant from the Research Grants Council of the Hong Kong Special Administrative Region, China HKUST C6004-25G and an ITF grant under the contract ITS/161/24FP. Dr. Yanan Guo and Dr. Zhenkai Zhang are supported by separate research gifts from HydroX AI, respectively.
\end{acks}

\bibliographystyle{ACM-Reference-Format}
\bibliography{ref}

\appendix
% \input{./files/appendix.tex}

\end{document}
